\section{Tree bounds and limits of the method}\label{app:limits}
\subsection{The input-occurrence charging lemma}\label{app:trees}
This subsection proves the tree refinement stated in
\S\ref{sec:trees}.

\treecharging*
\begin{proof}
For a subterm whose value on the reference structure is $X$, let
$S_X\subseteq\Lambda$ be the off-diagonal support of $X$. Boolean operations
act on these sets in the usual way and converse fixes them. By the
identities~\eqref{eq:table-off}, a composition $X;Y$ with both supports nonempty
has off-diagonal support $\Lambda$. Call such a node \defn{active}.
If either support is empty, the off-diagonal table of the product is
zero or one of the two operand tables, according to the diagonal bits.
Such an inactive node notices no fold: every outside-witness
contribution in the local calculation vanishes.

Construct a Boolean formula for the off-diagonal table of any subterm
by replacing active compositions by the constant true, inactive
compositions by zero or the relevant operand table, and converse by
its argument. This operation never duplicates an input occurrence.
For each active node $g=(X;Y)$, apply this pruning to its two operands.
Let $\mathcal L_g$ be the set of original input occurrences that remain
in the two resulting table formulas. These are the occurrences that
\emph{survive the pruning}, not all input occurrences below $g$.
This distinction prevents repeated charging at successive active nodes.
The sets $\mathcal L_g$ are
pairwise disjoint: an active descendant is replaced by a constant when
an ancestor's operands are pruned, and incomparable nodes have disjoint
subtrees.

By Lemma~\ref{lem:local}, a composition notices the fold of a color
$\lambda$ only if
\[
 S_X=\{\lambda\},\qquad S_Y=\{\lambda\},\qquad\text{or}\qquad
 S_X\cap S_Y=\{\lambda\}.
\]
Let $E_g$ be the set of colors satisfying at least one of these
conditions. Thus $|D_g|\le|E_g|$. If $S_X\cap S_Y$ is nonempty,
$E_g$ has size at most one. Its unique color, if present, is selected
by one operand table or by their conjunction. The selector assumption
therefore gives $|\mathcal L_g|\ge d|E_g|$. If the two supports are
disjoint, only singleton operands contribute to $E_g$. Each contributing
operand has at least $d$ input occurrences, and the operands are disjoint
trees. The same inequality follows. Summing over active nodes proves
\[
 \tcost(T)=\sum_g|D_g|\le\sum_g|E_g|
       \le\frac1d\sum_g|\mathcal L_g|\le \relatoms{T}/d.
\]
Diagonal bits may affect which operand table an inactive node retains, but do not
create additional input occurrences in this argument.
\end{proof}

\subsection{Factored trees for \texorpdfstring{$G_k$}{Gk}}\label{app:factored}
A complete binary decision tree successively chooses $P_i$ or $Q_i$ for
the read edge. At leaf $A$ place $H_{A^c};\one$. A branch at coordinate
$i$ with child terms $U_1,U_0$ has the form
\[
 (P_i\cap U_1)\cup(Q_i\cap U_0).
\]
Distributing the branches recovers the DNF for $G_k$ on arbitrary inputs;
no assumption that exactly one of $P_i,Q_i$ holds at each edge is used. Each of the $2^k$ leaves costs $2k+1$
nodes, and each of the $2^k-1$ internal branches adds five. The total is
$(2k+6)2^k-5$. The number of relation-input occurrences is
$k2^k+2(2^k-1)$.

\subsection{Why the requirements are crossed}
For comparison, the superficially similar non-crossed formula collapses:
\[
 \exists z\,\bigwedge_{i=1}^k
 [(P_i(x,y)\land P_i(x,z))\lor(Q_i(x,y)\land Q_i(x,z))]
 \equiv\bigwedge_{i=1}^k(P_i(x,y)\lor Q_i(x,y)).
\]
The forward implication is immediate; for the reverse implication take
$z=y$. In $G_k$, however, the two edges are subject to crossed
requirements, so choosing $z=y$ need not satisfy them.

\subsection{Sharpness of the fold tests}\label{app:promise}
\begin{proposition}[Sharpness on the folded structures]\label{prop:promise}
For every $k\ge1$, the minimum number of composition gates in a
relational circuit that computes $G_k$ on the reference structure and on
all its folded structures is $2^{k-1}$.
\end{proposition}
\begin{proof}
For each unordered complementary pair $\{A,A^c\}$ of colors, choose one
representative $A$ and compute, once,
\[
 X_A\coloneqq H_A;H_{A^c}.
\]
Use the shared result in
\[
 \bigcup_{\{A,A^c\}}
 \left[(H_{A^c}\cap X_A)\cup(H_A\cap X_A^\smile)\right].
\]
This circuit has exactly $2^{k-1}$ compositions. We check that it
computes $G_k$ on $\mathcal M_0$ and on every $\mathcal M_\lambda$. Let
$J$ be the set of colors at the hub, so that $J=\Lambda$ on
$\mathcal M_0$ and $J=\Lambda\setminus\{\lambda\}$ on
$\mathcal M_\lambda$. On each of these structures, $H_A$ is the set of
edges of color $A$. The relation $X_A$ is true on all pairs of distinct
non-hub vertices, has zero diagonal, and has hub-row and hub-column
values $[A\in J]$ and $[A^c\in J]$, respectively. The masks select the
complementary read colors. At a non-hub vertex every color still occurs,
so $G_k$ is true on every off-diagonal read edge there. At the hub it is
true precisely when the complementary color lies in $J$. These are
exactly the displayed circuit's values. The lower bound is
Theorem~\ref{thm:main}, whose proof uses only these structures.
\end{proof}
